\section[Classif\/ication of irreducible generic Gelfand--Tsetlin $\mathfrak{gl}(n)$-modules]{Classif\/ication of irreducible generic\\
Gelfand--Tsetlin $\boldsymbol{\mathfrak{gl}(n)}$-modules}
\label{section6}

In this section we prove the main result in the paper, i.e.~the generalization of Theorem~\ref{Basis for irreducible
generic modules gl(3)} for~$\mathfrak{gl}(n)$.
For convenience we introduce and recall some notation.

\begin{Notation}
\label{notation}
Let $T(L)=T(l_{ij})$ be a~f\/ixed tableau of height~$n$.
\begin{itemize}\itemsep=0pt
\item[(i)] $\mathcal{B}(T(L)):=\big\{T(L+z): z\in \mathbb{Z}^{\frac{n(n-1)}{2}}\big\}$.
\item[(ii)] $V(T(L)):= \Span \mathcal{B}(T(L))$.
\item[(iii)] For any $T(R)=T(r_{ij})\in \mathcal{B}(T(L))$ and for any $1<p\leq n$, $1\leq s\leq p$ and $1\leq u\leq
p-1$ we def\/ine:
\begin{itemize}\itemsep=0pt
\item[(a)] $\omega_{p,s,u}(T(R)):=r_{p,s}-r_{p-1,u}$;
\item[(b)] $\Omega(T(R)):=\{(p,s,u): \omega_{p,s,u}(T(R))\in \mathbb{Z}\}$;
\item[(c)] $\Omega^{+}(T(R)):=\{(p,s,u): \omega_{p,s,u}(T(R))\in \mathbb{Z}_{\geq 0}\}$;
\item[(d)] $\mathcal{N}(T(R)):=\{T(Q)\in \mathcal{B}(T(L)): \Omega^{+}(T(R))\subseteq \Omega^{+}(T(Q))\}$;
\item[(e)] $ W(T(R)):=\Span \mathcal{N}(T(R))$;
\item[(f)] $U\cdot T(R)$: the $\mathfrak{gl}(n)$-submodule of $V(T(L))$ generated by $T(R)$.
\end{itemize}
\end{itemize}
\end{Notation}

\subsection{Basis for the module generated by a~single tableau}
In order to f\/ind an explicit basis of every irreducible generic module, we f\/irst f\/ind a~basis of $U\cdot T(R)$ for any
tableau $T(R)$ in $\mathcal{B}(T(L))$.
\begin{Proposition}
\label{submodules of V(T(L))}
For any $T(R)\in\mathcal{B}(T(L))$, the Gelfand--Tsetlin formulas endow $ W(T(R))$ with a~$\mathfrak{gl}(n)$-module
structure.
\end{Proposition}
\begin{proof}
It is enough to prove $U\cdot T(Q)\subseteq W(T(R))$ for any $T(Q)=T(q_{ij})\in\mathcal{N}(T(R))$.
We will show $g\cdot T(Q)$ is in $W(T(R))$ for every (standard) generator~$g$ of $\mathfrak{gl}(n)$.

Suppose $g=E_{k,k+1}$ for some $1\leq k\leq n-1$.
By the Gelfand--Tsetlin formulas, we have
\begin{gather*}
E_{k,k+1}(T(Q))=-\sum\limits_{i=1}^{k}\left(\frac{\prod\limits_{j=1}^{k+1}(q_{ki}-q_{k+1,j})}{\prod\limits_{j\neq
i}^{k}(q_{ki}-q_{kj})}\right)T\big(Q+\delta^{ki}\big).
\end{gather*}
If $E_{k,k+1}(T(Q))\notin W(T(R))$, then there exist~$k$ and~$i$ such that $T(Q)\in \mathcal{N}(T(R))$ but
$T(Q+\delta^{ki})$ $\notin\mathcal{N}(T(R))$.
That implies
\begin{gather*}
\Omega^{+}(T(R))\subseteq\Omega^{+}(T(Q))~\text{and}~\Omega^{+}(T(R))\nsubseteq\Omega^{+}\big(T\big(Q+\delta^{ki}\big)\big).
\end{gather*}
Hence, there exists $(p,s,u)\in\Omega^{+}(T(R))$ such that $\omega_{p,s,u}(T(Q))\in\mathbb{Z}_{\geq 0}$ and
$\omega_{p,s,u}(T(Q+\delta^{ki}))$ $\notin\mathbb{Z}_{\geq 0}$.
The latter holds only in two cases:
\begin{gather*}
(p,s,u)\in \{(k,i,u), (k+1,s,i): 1\leq u\leq k-1; \, 1\leq s\leq k+1\}.
\end{gather*}
Note that if neither of these two cases hold, we have $\omega_{p,s,u}(T(Q+\delta^{ki}))=\omega_{p,s,u}(T(Q))$.
We consider now each of the two cases separately.
\begin{enumerate}[(i)]\itemsep=0pt
\item Suppose $(p,s,u)=(k,i,u)$.
Then $\omega_{k,i,u}(T(Q))=q_{ki}-q_{k-1,u}\in\mathbb{Z}_{\geq 0}$ and
$\omega_{k,i,u}(T(Q+\delta^{ki}))=(q_{ki}+1)-q_{k-1,u}\notin \mathbb{Z}_{\geq 0}$, which is impossible.
\item Suppose $(p,s,u)=(k+1,s,i)$.
Then
\begin{gather*}
\omega_{k+1,s,i}(T(Q))=q_{k+1,s}-q_{ki}\in\mathbb{Z}_{\geq 0}
\end{gather*}
{and}
\begin{gather*}
\omega_{k+1,s,i}(T(Q+\delta^{ki}))=q_{k+1,s}-(q_{ki}+1)\notin \mathbb{Z}_{\geq 0}.
\end{gather*}
Hence $q_{k+1,s}-q_{k,i}=0$ and then the coef\/f\/icient of $T(Q+\delta^{ki})$ in the decomposition of $E_{k,k+1}(T(Q))$ is
$-\frac{\prod\limits_{j=1}^{k+1}(q_{ki}-q_{k+1,j})}{\prod\limits_{j\neq i}^{k}(q_{ki}-q_{kj})}=0$.
\end{enumerate}
Therefore, the tableaux that appear with nonzero coef\/f\/icients in $E_{k,k+1}(T(Q))$ are elements of $N(T(R))$.
Hence, $E_{k,k+1}(T(Q))\in W(T(R))$.
The proof that $E_{k+1,k}(T(Q))\in W(T(R))$ is analogous to the one of $E_{k,k+1}(T(Q))$ $\in W(T(R))$.
The case $g=E_{kk}$ is trivial because $E_{kk}$ acts as a~multiplication by a~scalar on $T(Q)$ and
$T(Q)\in\mathcal{N}(T(R))\subseteq W(T(R))$.
\end{proof}

Given any tableau $T(R)$, there are three modules containing $T(R)$: $V(T(L))$, $ W(T(R))$ and $U\cdot T(R)$.
We will show that $ W(T(R))=U\cdot T(R)$.
For this we need the following lemmas.

\begin{Lemma}
\label{paths for connected tableaux}
Let $T(L)$ be a~generic tableau.
If $0\neq z\in\mathbb{Z}^{\frac{n(n-1)}{2}}$ is such that $\Omega^{+}(T(L))\subseteq\Omega^{+}(T(L+z))$ then, there
exist $i$, $j$ such that $z_{ij}\neq 0$ and
\begin{gather}
\label{inclutions of Omega +}
\Omega^{+}(T(L))\subseteq\Omega^{+}\big(T\big(L+z_{ij}\delta^{ij}\big)\big)\subseteq\Omega^{+}(T(L+z)).
\end{gather}
\end{Lemma}
\begin{proof}
We will use the following def\/inition in the proof of the lemma.

\begin{Definition}
Given a~generic tableau $T(R)\in\mathcal{B}(T(L))$, a~\emph{chain in $T(R)$ of length~$\ell$ starting in row~$d$} is
a~subset of the entries of $T(R)$, $C=\{r_{d-i,s^{(d-i)}}\}_{i=0,\ldots, \ell}$, where $1\leq s^{(d-i)}\leq d-i$ are
such that $r_{d-i,s^{(d-i)}}-r_{d-i-1,s^{(d-i-1)}}\in\mathbb{Z}$ for any $i=0,\ldots,\ell-1$
(i.e.~$\{(d-i,s^{(d-i)},s^{(d-i-1)})\}_{i=0,\ldots,\ell}$ $\subseteq\Omega(T(R))$). The chain is called maximal if
\begin{itemize}\itemsep=0pt
\item[(i)] $(d+1,i,s^{(d)})\notin\Omega(T(R))$ for any $1\leq i\leq d+1$,
\item[(ii)] $(d-\ell,s^{(d-\ell)},j)\notin\Omega(T(R))$ for any $1\leq j\leq d-\ell-1$.
\end{itemize}
\end{Definition}

For every $T(R)$ in $\mathcal{B}(T(L))$ we have that $\Omega^{+}(T(R))= \bigsqcup_{1\leq c\leq n}\Omega^{+}_{c}(T(R))$,
where $\Omega^{+}_{c}(T(R)):=\{(p,s,u)\in\Omega^{+}(T(R)): p=c\}$.
In particular,~\eqref{inclutions of Omega +} holds if and only if
\begin{gather}
\label{inclutions for Omega c +}
\Omega^{+}_{c}(T(L))\subseteq\Omega^{+}_{c}\big(T\big(L+z_{ij}\delta^{ij}\big)\big)\subseteq\Omega^{+}_{c}(T(L+z))
\end{gather}
for any $1\leq c\leq n$.
For $c\notin\{i,i+1 \}$ we have $\Omega^{+}_{c}(T(L))=\Omega^{+}_{c}(T(L+z_{ij}\delta^{ij}))$.
So, in order to verify~\eqref{inclutions for Omega c +}, it is enough to consider the cases $c=i, i+1$.

Lets consider $k$, $l$ such that $z_{kl}\neq 0$.
Set for convenience $Q:=L+z$.
There exists a~maximal chain~$C$ in $T(Q)$ of length~$\ell$, starting in row~$d$ such that $q_{kl}\in C$.
Suppose that $C=\{q_{[i]}\}_{i=0,\ldots,\ell}$ where $[i]:=(d-i,s^{(d-i)})$.
If $\ell=0$, then $C=\{q_{kl}\}$ and~\eqref{inclutions of Omega +} is obvious for $z_{ij}=z_{kl}$.

Let~$a$ and~$b$ be the minimum and maximum of $\{i:z_{[i]}\neq 0\}$, respectively.
We have
\begin{gather}
\Omega^{+}_{d-a+1}\big(T\big(L+z_{[a]}\delta^{[a]}\big)\big)=\Omega^{+}_{d-a+1}(T(L+z)),\nonumber
\\
\Omega^{+}_{d-b}\big(T\big(L+z_{[b]}\delta^{[b]}\big)\big)=\Omega^{+}_{d-b}(T(L+z)).\label{a and b behavior}
\end{gather}

Therefore~\eqref{inclutions for Omega c +} holds for the pairs $c=d-a+1$, $z_{ij}=z_{[a]}$ and $c=d-b$, $z_{ij}=z_{[b]}$,
respectively.
Now, let $a\leq m\leq b$ and consider the $4$ cases depending on what the signs of $z_{[a]}$ and $z_{[a+1]}$ are.
\begin{enumerate}[(i)]\itemsep=0pt
\item $z_{[m]}>0$ and $z_{[m+1]}\leq 0$.
In this case~\eqref{inclutions for Omega c +} holds for $c=d-m$ and $z_{ij}=z_{[m]}$.
In particular, if $z_{[a]}>0$ and $z_{[a+1]}\leq 0$, using the f\/irst equation in~\eqref{a and b behavior}, we conclude
that~\eqref{inclutions of Omega +} holds for $z_{ij}=z_{[a]}$.
\item $z_{[m]}<0$ and $z_{[m-1]}\geq 0$.
In this case~\eqref{inclutions for Omega c +} holds for $c=d-m+1$ and $z_{ij}=z_{[m-1]}$.
In particular, if $z_{[b]}<0$ and $z_{[b-1]}\geq 0$, using the second equation in~\eqref{a and b behavior} we conclude
that~\eqref{inclutions of Omega +} holds for $z_{ij}=z_{[b]}$.
\item $z_{[m]}>0$ and $z_{[m+1]}> 0$.
In this case~\eqref{inclutions for Omega c +} holds for $c=d-m$ and
\begin{gather*}
z_{ij}=
\begin{cases}
z_{[m]} & \text{if} \ \ l_{[m]}-l_{[m+1]}\in\mathbb{Z}_{\geq 0},
\\
z_{[m+1]} & \text{if} \ \ l_{[m+1]}-l_{[m]}\in\mathbb{Z}_{> 0}.
\end{cases}
\end{gather*}
\item $z_{[m]}<0$ and $z_{[m-1]}< 0$.
In this case~\eqref{inclutions for Omega c +} holds for $c=d-m+1$ and
\begin{gather*}
z_{ij}=
\begin{cases}
z_{[m]} & \text{if} \ \  l_{[m-1]}-l_{[m]}\in\mathbb{Z}_{\geq 0},
\\
z_{[m-1]} & \text{if} \ \ l_{[m]}-l_{[m-1]}\in\mathbb{Z}_{> 0}.
\end{cases}
\end{gather*}
\end{enumerate}

Now combining (i)--(iv) we reduce the proof to the following two cases:
\begin{enumerate}[(a)]\itemsep=0pt
\item $z_{[a]}>0, z_{[a+1]}>0, \ldots, z_{[b]}>0$ and for any $t=1,\ldots,b-a$,~\eqref{inclutions for Omega c +} holds
for $c=d-a+t+1$ and $z_{ij}=z_{[a+l]}$.
In particular,~\eqref{inclutions for Omega c +} holds for $c=d-b+1$ and $z_{ij}=z_{[b]}$.
So, by the second equation in~\eqref{a and b behavior} we have that~\eqref{inclutions of Omega +} holds for
$z_{ij}=z_{[b]}$.
\item $z_{[b]}<0, z_{[b-1]}<0, \ldots, z_{[a]}<0$ and for any $t=1,\ldots,b-a$,~\eqref{inclutions for Omega c +} holds
for $c=d-(b-t)$ and $z_{ij}=z_{[b-t]}$.
In particular,~\eqref{inclutions for Omega c +} holds for $c=d-a$ and $z_{ij}=z_{[a]}$.
So, by the f\/irst equation in~\eqref{a and b behavior} we have that~\eqref{inclutions of Omega +} holds for $z_{ij}=z_{[a]}$. \hfill{$\qed$}
\end{enumerate}
\renewcommand{\qed}{}
\end{proof}


\begin{Definition}
Given $T(Q)$ and $T(R)$ in $\mathcal{B}(T(L))$, we write $T(R)\preceq_{(1)}T(Q)$ if there exist $g\in\mathfrak{gl}(n)$
such that $T(Q)$ appears with nonzero coef\/f\/icient in the decomposition of $g\cdot T(R)$ into a~linear combination of
tableaux.
For any $p\geq 1$ we write $T(R)\preceq_{(p)} T(Q)$ if there exist tableaux $T(L^{(1)})$,\dots, $T(L^{(p)})$, such that
\begin{gather*}
T(R)=T\big(L^{(0)}\big)\preceq_{(1)}T\big(L^{(1)}\big)\preceq_{(1)}\cdots\preceq_{(1)}T\big(L^{(p)}\big)=T(Q).
\end{gather*}
\end{Definition}

As an immediate consequence of the def\/inition of $\preceq_{(p)}$ we have the following.
\begin{Lemma}
\label{properties of relation less or equal}
If $T(Q)$, $T(Q^{(0)})$, $T(Q^{(1)})$ and $T(Q^{(2)})$ are tableaux in $\mathcal{B}(T(L))$ then:
\begin{itemize}\itemsep=0pt
\item[$(i)$] $T(Q^{(0)})\preceq_{(p)} T(Q^{(1)})$ and $T(Q^{(1)})\preceq_{(q)} T(Q^{(2)})$ imply
$T(Q^{(0)})\preceq_{(p+q)} T(Q^{(2)})$;
\item[$(ii)$] $T(Q)\preceq_{(1)} T(Q)$.
\end{itemize}
\end{Lemma}

\begin{Corollary}\label{Gelfand--Tsetlin modules property}\sloppy
If $T(R), T(Q)\in \mathcal{B}(T(L))$ are generic Gelfand--Tsetlin tableaux such that
\mbox{$T(R)\preceq_{(p)} T(Q)$} for some $p\in\mathbb{Z}_{\geq 0}$, then $T(Q)\in U\cdot T(R)$.
\end{Corollary}
\begin{proof}
By Lemma~\ref{elements of Gamma separates tableaux} and the def\/inition of the relation $\preceq_{(1)}$, we f\/irst verify
that ${T(R)\preceq_{(1)}\! T(Q)}$ implies $T(Q)\in U\cdot T(R)$.
Now, using Lemma~\ref{properties of relation less or equal}(i), if $T(R)\preceq_{(p)} T(Q)$ for some~$p$ then $T(Q)\in
U\cdot T(R)$.
\end{proof}

The next theorem provides a~convenient basis for the submodule of $V(T(L))$ generated by a~f\/ixed tableau.
Recall the def\/inition of $\mathcal{N}(T(R))$ in Notation~\ref{notation}(iii)(d).

\begin{Theorem}
\label{basis of submodule generated by T}
For any tableau $T(R)\in\mathcal{B}(T(L))$, $U\cdot T(R)= W(T(R))$.
In particular, $\mathcal{N}(T(R))$ forms a~basis of $U\cdot T(R)$, and the action of $\mathfrak{gl}(n)$ on $U\cdot T(R)$
is given by the Gelfand--Tsetlin formulas.
\end{Theorem}

\begin{proof}
By Proposition~\ref{submodules of V(T(L))}, $U\cdot T(R)\subseteq W(T(R))$.
To prove that $ W(T(R))\subseteq U\cdot T(R)$ we will show that $T(Q)\in U\cdot T(R)$ for any
$T(Q)\in\mathcal{N}(T(R))$.
By Corollary~\ref{Gelfand--Tsetlin modules property}, it is enough to prove that $T(R)\preceq_{(p)} T(Q)$ for some
positive integer~$p$.

Suppose that $T(Q)=T(R+z)\in \mathcal{N}(T(R))$ for some $z\in\mathbb{\mathbb{Z}}^{\frac{n(n-1)}{2}}$.
Let~$t$ be the number of non-zero components of~$z$.
We will prove that $T(R)\preceq_{(p)} T(Q)$ using induction on~$t$.

Let us f\/irst consider the case $t=1$ (the case $t=0$ is trivial, since then $T(Q)=T(R)$) and $z_{ij}>0$.
We will f\/irst prove that $T(R+l\delta^{ij})\preceq_{(1)}T(R+(l+1)\delta^{ij})$ for any $0\leq l \leq z_{ij}-1$.
This will imply
\begin{gather*}
T(R)\preceq_{(1)}T\big(R+\delta^{ij}\big)\preceq_{(1)}T\big(R+2\delta^{ij}\big)\preceq_{(1)}\cdots\preceq_{(1)}T\big(R+z_{ij}\delta^{ij}\big)=T(Q),
\end{gather*}
and then $T(R)\preceq_{(z_{ij})}T(Q)$.
To prove that $T(R+l\delta^{ij})\preceq_{(1)}T(R+(l+1)\delta^{ij})$ we show that the coef\/f\/icient of
$T(R+(l+1)\delta^{ij})$ in the decomposition of $E_{i,i+1}(T(R+l\delta^{ij}))$ is not zero.
In fact, by the Gelfand--Tsetlin formulas, that coef\/f\/icient is
\begin{gather*}
a_{l}:=-\frac{\prod\limits_{k=1}^{i+1}(r_{ij}-r_{i+1,k}+l)}{\prod\limits_{k\neq j}^{i}(r_{ij}-r_{ik}+l)}.
\end{gather*}
Assume that $a_{l}=0$.
Then $r_{ij}-r_{i+1,k}+l=0$ for some~$k$, which implies $\omega_{i+1,k,j}(T(R))$ $=r_{i+1,k}-r_{ij}=l\in\mathbb{Z}_{\geq
0}$.
But, since $T(Q)\in\mathcal{N}(T(R))$, we have{\samepage
\begin{gather*}
l-z_{ij}=r_{i+1,k}-r_{ij}-z_{ij}=\omega_{i+1,k,j}(T(Q))\in\mathbb{Z}_{\geq{0}}.
\end{gather*}
Therefore we have %both
$ 0\leq l \leq z_{ij}-1$ and $z_{ij}\leq l$, which is a~contradiction.
Hence,
$T(R)\preceq_{(z_{ij})}T(Q)$.}

Let now $t=1$ and $z_{ij}<0$.
Using the same arguments as in the case $z_{ij}>0$, we prove that $T(R)\preceq_{(-z_{ij})}T(Q)$ using $|z_{ij}|$
applications of $E_{i+1,i}$.
This completes the proof for $t=1$.

Assume now that for any $w\in\mathbb{Z}^{\frac{n(n-1)}{2}}$ with at most~$t$ nonzero components, and such that
$\Omega^{+}(T(R))\subseteq\Omega^{+}(T(R+w))$, we have $T(R)\preceq_{(p)} T(R+w)$ for some~$p$.
Let us consider~$z$ with $t+1$ nonzero components.
Since $\Omega^{+}(T(R))\subseteq \Omega^{+}(T(R+z))$, by Lemma~\ref{paths for connected tableaux}, there exist $i$, $j$
such that
\begin{gather*}
\Omega^{+}(T(R))\subseteq\Omega^{+}\big(T\big(R+z_{ij}\delta^{ij}\big)\big)\subseteq \Omega^{+}(T(R+z)).
\end{gather*}
Using the induction hypothesis for the pairs of tableaux $(T(R), T(R+z_{ij}\delta^{ij}))$ and
$(T(R+z_{ij}\delta^{ij}),T(R+z))$, there exist $p,q\in\mathbb{Z}_{\geq 0}$ such that
\begin{gather*}
T(R)\preceq_{(p)} T\big(R+z_{ij}\delta^{ij}\big)\qquad \text{and} \qquad T\big(R+z_{ij}\delta^{ij}\big)\preceq_{(q)} T(R+z).
\end{gather*}
Thus, by Lemma~\ref{properties of relation less or equal}(i), $T(R)\preceq_{(p+q)}T(R+z)$.
\end{proof}

\begin{Proposition}
\label{when two tableaux generate the same submodule}
Let $T(R)$ and $T(Q)$ be in $\mathcal{B}(T(L))$.
Then $U\cdot T(R)= U\cdot T(Q)$ if and only if $\Omega^{+}(T(Q))=\Omega^{+}(T(R))$.
\end{Proposition}
\begin{proof}\sloppy
Using Theorem~\ref{basis of submodule generated by T} and the def\/initions of $ W(T(R))$, $W(T(Q))$, $\Omega^{+}(T(R))$,
and $\Omega^{+}(T(Q))$, we can prove a~stronger statement: $U\cdot T(R) \subseteq U\cdot T(Q)$ if and only if
$\Omega^{+}(T(Q))\subseteq\Omega^{+}(T(R))$.
\end{proof}

\begin{Corollary}
$U \cdot T(R)=V(T(L))$ whenever $\Omega^{+}(T(R))=\varnothing$.
\end{Corollary}

\begin{Definition}
We will write $T(Q)\sim_{\Omega^+}T(R)$ if $\Omega^{+}(T(R)) = \Omega^{+}(T(Q))$.
\end{Definition}

\begin{Proposition}
%\label{every submodule of V(T(L)) is fin gen}
Every submodule of $V(T(L))$ is finitely generated.
\end{Proposition}
\begin{proof}
Let~$N$ be any submodule of $V(T(L))$ and~$\Phi$ the set of all tableaux $T(R)$ in~$N$ such that
$\Omega^{+}(T(P))\subseteq\Omega^{+}(T(R))$ implies $\Omega^{+}(T(P))=\Omega^{+}(T(R))$.
By Theorem~\ref{basis of submodule generated by T}, $N=\sum\limits_{T(R)\in\Phi} U\cdot T(R)$ and by
Proposition~\ref{when two tableaux generate the same submodule}, we can write $N=\bigoplus_{T(R)\in\tilde{\Phi}} U\cdot
T(R)$, where $\tilde{\Phi}$ is a~set of distinct representatives of $\Phi/\sim_{\Omega^+}$ (hence
$\Omega^{+}(T(R))\neq\Omega^{+}(T(Q))$ for any $T(R)$, $T(Q)$ in $\tilde{\Phi}$).
Now, since $\Omega(T(L))$ is a~f\/inite set, then $\tilde{\Phi}$ is f\/inite.
\end{proof}

\subsection{Basis for irreducible modules containing a~given tableau}
By Theorem~\ref{basis of submodule generated by T}, the module generated by a~tableau $T(R)$ has basis
$\mathcal{N}(T(R))$.
